\section{Conclusion}
\label{sec:prof-conclusion}

ORIUS argues that the right universal object for Physical AI safety is not a single
controller, a single benchmark, or a single proof family. It is a runtime safety layer
that mediates degraded observation before actuation. Once telemetry reliability becomes
part of the safety problem, the system must expose reliability, uncertainty inflation,
tightened action legality, repair, fallback, and certificate semantics as first-class
objects. That is the contribution of the Detect--Calibrate--Constrain--Shield--Certify
kernel.

The manuscript makes that contribution at three levels. First, it defines OASG as the
dominant hidden hazard: actions can remain legal on the observed state while becoming
unsafe on the true state. Second, it shows how a universal runtime contract can bind
adapter semantics, benchmark semantics, and CertOS governance under one typed layer.
Third, it demonstrates that the architecture is more than a diagram: the battery row
provides witness-depth theorem-to-code-to-artifact closure, and the defended bounded
rows show that the same runtime grammar travels beyond energy.

The scientific boundary is intentionally explicit. The paper does not claim equal-depth
closure beyond the defended three-domain lane. Battery is the witness row, and AV plus
Healthcare are bounded defended rows. Future additional domains remain outside the
current defended evidence package. That boundary does not weaken the architectural
thesis. It shows that ORIUS is disciplined enough to separate universality of contract
from the present evidence depth.

What changes at the field level is the object of research. ORIUS suggests that Physical
AI safety should be organized around reusable runtime layers that sit between nominal
intelligence and physical release. Such a layer can host conformal uncertainty, safety
filters, runtime assurance, and anomaly detection without collapsing into any one of
them. If future additional domains clear the same obligations, ORIUS supports a
broader empirical claim. Even before that point, it establishes a credible
safety-layer grammar for Physical AI under degraded observation.

In that sense, ORIUS is also a writing and evaluation proposal for the field. It argues
that future physical-AI systems should report not only predictive accuracy or planner
reward, but also true-state violation behavior, intervention frequency, fallback usage,
certificate validity, and release latency under degraded observation. That reporting
discipline matters because many failures in real systems are not failures of nominal
optimization; they are failures of release semantics once observation ceases to be
trustworthy. A framework like ORIUS makes those failures measurable and governable.
